\begin{frame}{리뷰 쓰는 과정 — 체크리스트 $\to$ 반박 가능한 질문}
  \begin{columns}[T]
    \begin{column}{0.5\textwidth}
      \begin{itemize}
        \item (1) 다른 팀 발표 + 보고서 읽고 \textbf{주장 파악}
          (알고리즘·하이퍼파라미터·숫자·결론)
        \item (2) \textbf{체크리스트 5항목} 적용
        \item (3) 각 항목에서 \textbf{질문 초안} — 답이 \emph{새
          계산·실험}이어야 함
        \item (4) \textbf{분기}: ``Ch9$\sim$15 개념(장·절)을 근거로
          들고, 답하려면 \textbf{새 실험}(시드 재실행·평가 프로토콜
          변경·베이스라인 실측·스케일 비교)이 필요한가?''
          \begin{itemize}
            \item 예 $\to$ \textbf{반박 가능한 질문 채택}(초록)
            \item 아니오 $\to$ ``곡선이 올라가네요'' 같은 \textbf{감상평}
              (빨강)으로 떨어짐 — Ch09.2($\varepsilon$-greedy)·Ch11.3
              (GAE $\lambda$)·Ch12.1(복합 오차) 개념으로 \textbf{다시
              쓰기}(점선으로 (3) 회귀)
          \end{itemize}
        \item (5) 3단계 기준(1/2/3점)으로 \textbf{자기 채점}
      \end{itemize}
    \end{column}
    \begin{column}{0.47\textwidth}
      \centering
      \includegraphics[width=\linewidth]{ch16_2_review_flow.png}
      \vskip0.3em
      {\footnotesize 신경망 RL 버전의 리뷰 플로우.}
    \end{column}
  \end{columns}
\end{frame}
